%! Diagonalizing a Matrix — Unit 6, Lesson 6.2 — Homework (Answer Key)
\documentclass[10pt]{article}
\usepackage{linalg-article}
\usepackage{linalg-key}   % pulls in -boxes; do NOT also load -boxes
\newcommand{\TallMath}[1]{$\displaystyle #1\rule[-1.4em]{0pt}{3.2em}$}
\newcommand{\xx}{\mathbf{x}}
\newcommand{\zero}{\mathbf{0}}

\begin{document}

\pageheader{Unit 6, Lesson 6.2}{Homework --- Answer Key}
\namedateperiod

\vspace{0.12in}

\begin{notesbox}{Practice}
Line the eigenvectors into the columns of $X$ and the eigenvalues down the diagonal of $\Lambda$. Then
$A=X\Lambda X^{-1}$, and $A^k=X\Lambda^k X^{-1}$.

\begin{enumerate}[leftmargin=1.9em, itemsep=12pt]
  \item \textbf{Diagonalize.} For $A=\begin{bmatrix} 1 & 2 \\ 2 & 1 \end{bmatrix}$ (eigenvalues $\lambda=3,-1$
        with eigenvectors $\begin{bsmallmatrix}1\\1\end{bsmallmatrix},\begin{bsmallmatrix}1\\-1\end{bsmallmatrix}$),
        write $X$, $\Lambda$, and $X^{-1}$, then write $A=X\Lambda X^{-1}$.
        \ansline{$X=\begin{bsmallmatrix}1&1\\1&-1\end{bsmallmatrix}$, $\Lambda=\begin{bsmallmatrix}3&0\\0&-1\end{bsmallmatrix}$; $\det X=-2$, so $X^{-1}=\begin{bsmallmatrix}\frac12&\frac12\\[2pt]\frac12&-\frac12\end{bsmallmatrix}$. Then $A=X\Lambda X^{-1}=\begin{bsmallmatrix}3&-1\\3&1\end{bsmallmatrix}\begin{bsmallmatrix}\frac12&\frac12\\[2pt]\frac12&-\frac12\end{bsmallmatrix}=\begin{bsmallmatrix}1&2\\2&1\end{bsmallmatrix}.\ \checkmark$}
  \item \textbf{Powers made easy.} For $A=\begin{bmatrix} 2 & 3 \\ 0 & 1 \end{bmatrix}$ (eigenvalues
        $\lambda=2,1$ with eigenvectors $\begin{bsmallmatrix}1\\0\end{bsmallmatrix},\begin{bsmallmatrix}3\\-1\end{bsmallmatrix}$),
        use $A^4=X\Lambda^4 X^{-1}$ to compute $A^4$ \emph{without} multiplying $A$ four times.
        \ansline{$X=\begin{bsmallmatrix}1&3\\0&-1\end{bsmallmatrix}$, $\det X=-1$, $X^{-1}=\begin{bsmallmatrix}1&3\\0&-1\end{bsmallmatrix}$; $\Lambda^4=\begin{bsmallmatrix}16&0\\0&1\end{bsmallmatrix}$. So $A^4=X\Lambda^4X^{-1}=\begin{bsmallmatrix}16&3\\0&-1\end{bsmallmatrix}\begin{bsmallmatrix}1&3\\0&-1\end{bsmallmatrix}=\begin{bsmallmatrix}16&45\\0&1\end{bsmallmatrix}$.}
  \item \textbf{Diagonalizable or not?} For each matrix, decide whether it can be diagonalized and say why.
    \begin{enumerate}[label=(\alph*), leftmargin=1.8em, itemsep=4pt]
      \item $A=\begin{bmatrix} 3 & 0 \\ 1 & 3 \end{bmatrix}$ \hfill \ans{No.}
      \item $A=\begin{bmatrix} 2 & 1 \\ 1 & 2 \end{bmatrix}$ \hfill \ans{Yes.}
    \end{enumerate}
    \ansline{(a) Repeated $\lambda=3$; $(A-3I)=\begin{bsmallmatrix}0&0\\1&0\end{bsmallmatrix}\Rightarrow x_1=0\Rightarrow$ only one eigenvector direction $\begin{bsmallmatrix}0\\1\end{bsmallmatrix}$, so $X$ can't be invertible --- \emph{not} diagonalizable. (b) Two \emph{different} eigenvalues $\lambda=3,1$ $\Rightarrow$ two independent eigenvectors $\Rightarrow$ diagonalizable.}
  \item \textbf{Reading powers off $\Lambda$.} A matrix has $A=X\Lambda X^{-1}$ with
        $\Lambda=\begin{bmatrix} 2 & 0 \\ 0 & 5 \end{bmatrix}$ (the eigenvectors don't change). Without
        computing $A$, state the eigenvalues of $A$, of $A^{3}$, and of $A^{-1}$.
        \ansline{$A$: $\lambda=2,5$. $A^3=X\Lambda^3X^{-1}$: $\lambda=8,125$. $A^{-1}=X\Lambda^{-1}X^{-1}$: $\lambda=\frac12,\frac15$. (Same eigenvectors throughout --- only the eigenvalues change.)}
  \item \textbf{Explain.} Why does the factorization collapse to $A^k=X\Lambda^k X^{-1}$, and why must $A$
        have \emph{independent} eigenvectors for this to work?
        \ansline{Writing $A^k=(X\Lambda X^{-1})\cdots(X\Lambda X^{-1})$, each inner $X^{-1}X$ collapses to $I$, leaving $X\Lambda^k X^{-1}$. The eigenvectors must be independent so that $X$ is invertible --- without $X^{-1}$ there is no factorization at all.}
\end{enumerate}
\end{notesbox}

\begin{extensionbox}
\textbf{Back to the notes matrix.} In the notes you found $A^2$ for $A=\begin{bmatrix} 3 & 1 \\ 0 & 2 \end{bmatrix}$
(with $X=\begin{bsmallmatrix}1&1\\0&-1\end{bsmallmatrix}$, $X^{-1}=\begin{bsmallmatrix}1&1\\0&-1\end{bsmallmatrix}$).
Now use $A^4=X\Lambda^4 X^{-1}$ to compute $A^4$.
\ansline{$\Lambda^4=\begin{bsmallmatrix}81&0\\0&16\end{bsmallmatrix}$, so $A^4=X\Lambda^4X^{-1}=\begin{bsmallmatrix}81&16\\0&-16\end{bsmallmatrix}\begin{bsmallmatrix}1&1\\0&-1\end{bsmallmatrix}=\begin{bsmallmatrix}81&65\\0&16\end{bsmallmatrix}$.}
\end{extensionbox}

\begin{spiralbox}
\textbf{Coming next --- Lesson 6.3: Symmetric Positive Definite Matrices.} You saw today that a matrix
diagonalizes only when it has enough independent eigenvectors. Next you'll meet a family that
\emph{always} does --- \textbf{symmetric} matrices --- and whose eigenvectors are even \emph{perpendicular},
giving the cleanest diagonalization of all.
\end{spiralbox}

\begin{teachernote}
Item~1 is a full diagonalization on a symmetric matrix ($\lambda=3,-1$; the $X^{-1}$ has halves --- a good
place to reinforce the Unit~2 formula). Item~2 is the power payoff with a \emph{fraction-free} $X$
($A^4=\begin{bsmallmatrix}16&45\\0&1\end{bsmallmatrix}$). Item~3 is the diagonalizability decision (repeated
eigenvalue, one eigenvector $\Rightarrow$ no; distinct eigenvalues $\Rightarrow$ yes). Item~4 checks that
powers and inverses only change the \emph{eigenvalues} ($\Lambda^k$, $\Lambda^{-1}$). Item~5 is the target
justification (the $X^{-1}X$ cancellation; $X$ must be invertible). The extension revisits the notes matrix
one power further ($A^4=\begin{bsmallmatrix}81&65\\0&16\end{bsmallmatrix}$). Most common slips: column vs.\
eigenvalue order in $X$/$\Lambda$; fraction errors in $X^{-1}$ (item~1); and treating a repeated eigenvalue
as automatically non-diagonalizable (it's the \emph{missing second eigenvector} that blocks it).
\end{teachernote}

\end{document}
